\section{10. Limitations}

\subsection{10.1 Encoder-Only Benefit}

Galerkin provides linear complexity only for \textbf{encoder} models. Decoders 
fall back to standard attention.

\subsection{10.2 Not a Compression Method}

Parameter count increases by ~{{PARAM\_CHANGE}}\% via LoRA.

\subsection{10.3 Fixed Basis}

The DCT basis is fixed, not learned. A data-adaptive basis might be better.

\subsection{10.4 Smoothness Assumption}

DCT assumes smoothness. Attention patterns with sharp peaks (e.g., exact 
token matching) may be poorly represented by low-rank DCT.

\subsection{10.5 No Formal Error Bound on Downstream Accuracy}

We bound the operator approximation error (Section 4.9), but not the 
downstream task accuracy.

\subsection{10.6 Task Coverage}

Only SST-2 at $N=128$. Generation, QA, long-context untested.

\subsection{10.7 Model Coverage}

Three models (66.9M--124.6M).

\subsection{10.8 Hardware}

Only RTX 4070 Laptop.

\subsection{10.9 Rank Selection}

Fixed at $r = 64$. Rank ablation shows sensitivity.

\subsection{10.10 Numerical Stability}

DCT orthonormalization uses QR with $10^{-4}$ regularization. Extreme 
sequences may still have issues.

\subsection{10.11 Causal Cross-Attention}

Not tested.

\subsection{10.12 Code Quality}

\begin{itemize}
\item Research-grade
\end{itemize}
\begin{itemize}
\item Broad exception handling in \texttt{inject\_galerkin} (returns None on any failure)
\end{itemize}
\begin{itemize}
\item No unit tests
\end{itemize}

\subsection{10.13 Ethical Considerations}

Dual-use. See Appendix E.